\documentclass[11pt,a4paper]{article}
% =====================================================================
%  Shared preamble for the Part V lecture notes (lectures 19-22).
%
%  These four lectures sit outside Geron, so the notes ARE the primary
%  source and are examined on the same terms as the textbook parts.
%  Everything here is chosen to make that readable on paper: one column,
%  generous measure, and the same six-colour palette as the slides so a
%  student moving between the two is not re-learning what red means.
%
%  Build with pdflatex (twice, for the toc and the refs):
%      cd notes && latexmk -pdf lecture-19.tex
%  or  make -C notes
% =====================================================================
\usepackage[T1]{fontenc}
\usepackage[utf8]{inputenc}
\usepackage{newtxtext,newtxmath}      % Times-like: prints well, reads well
\usepackage[scaled=0.86]{inconsolata} % code, at a size that sits beside Times
\usepackage{microtype}
% newtx defines \Bbbk, \openbox, \proof and \endproof; amssymb and amsthm
% then redefine them and abort. Freeing the four names before those two
% packages load is the standard fix, and it costs nothing here: we use
% amsthm only for \newtheorem.
\let\Bbbk\relax
\usepackage{amsmath,amssymb}
\let\openbox\relax
\let\proof\relax
\let\endproof\relax
\usepackage{amsthm}
\usepackage{booktabs}
\usepackage{enumitem}
\usepackage{xcolor}
\usepackage{tcolorbox}
\usepackage{titlesec}
\usepackage{graphicx}
\usepackage[hidelinks]{hyperref}
\tcbuselibrary{skins,breakable}

% ---- the course palette, from assets/css/custom.css -------------------
\definecolor{aimlprimary}{HTML}{0B3D62}   % deep blue  - structure
\definecolor{aimlaccent} {HTML}{C0392B}   % brick red  - failures
\definecolor{aimlsuccess}{HTML}{14663A}   % green      - repairs
\definecolor{aimlmath}   {HTML}{6C3483}   % purple     - the derivation
\definecolor{aimlmuted}  {HTML}{4B5563}
\definecolor{aimlrule}   {HTML}{B0BCC7}
\definecolor{aimlsoft}   {HTML}{F4F7F9}

\usepackage[a4paper,margin=32mm,top=30mm,bottom=30mm]{geometry}
\setlength{\parskip}{0.45\baselineskip}
\setlength{\parindent}{0pt}
\linespread{1.06}

% ---- headings --------------------------------------------------------
\titleformat{\section}{\Large\bfseries\color{aimlprimary}}{\thesection}{0.7em}{}
\titleformat{\subsection}{\large\bfseries\color{aimlprimary}}{\thesubsection}{0.7em}{}
\titleformat{\subsubsection}{\bfseries\color{aimlmuted}}{}{0em}{}
\titlespacing*{\section}{0pt}{1.6\baselineskip}{0.5\baselineskip}
\titlespacing*{\subsection}{0pt}{1.2\baselineskip}{0.35\baselineskip}

% ---- boxes -----------------------------------------------------------
% One box per role, and the role is the colour. A student who has seen the
% slides already knows what each one means before reading a word of it.
\newtcolorbox{keybox}[1][]{
  breakable, enhanced, colback=aimlsoft, colframe=aimlprimary,
  boxrule=0.4pt, left=8pt, right=8pt, top=6pt, bottom=6pt,
  fonttitle=\bfseries\small, coltitle=white, colbacktitle=aimlprimary,
  attach boxed title to top left={yshift=-2mm, xshift=4mm},
  boxed title style={boxrule=0pt, arc=1pt}, #1}

\newtcolorbox{failbox}[1][]{
  breakable, enhanced, colback=white, colframe=aimlaccent,
  boxrule=0.9pt, left=8pt, right=8pt, top=6pt, bottom=6pt,
  fonttitle=\bfseries\small, coltitle=white, colbacktitle=aimlaccent,
  attach boxed title to top left={yshift=-2mm, xshift=4mm},
  boxed title style={boxrule=0pt, arc=1pt}, #1}

\newtcolorbox{derivbox}[1][]{
  breakable, enhanced, colback=white, colframe=aimlmath,
  boxrule=0.6pt, left=8pt, right=8pt, top=6pt, bottom=6pt,
  fonttitle=\bfseries\small, coltitle=white, colbacktitle=aimlmath,
  attach boxed title to top left={yshift=-2mm, xshift=4mm},
  boxed title style={boxrule=0pt, arc=1pt}, #1}

% An exercise. Deliberately the same shape as the notebook's `try` field:
% one change, and what should happen to the output.
\newtcolorbox{trybox}[1][]{
  breakable, enhanced, colback=white, colframe=aimlsuccess,
  boxrule=0.5pt, left=8pt, right=8pt, top=6pt, bottom=6pt,
  fonttitle=\bfseries\small, coltitle=white, colbacktitle=aimlsuccess,
  attach boxed title to top left={yshift=-2mm, xshift=4mm},
  boxed title style={boxrule=0pt, arc=1pt}, #1}

% ---- small semantics -------------------------------------------------
\newcommand{\scope}[1]{\textcolor{aimlmuted}{\small #1}}
\newcommand{\term}[1]{\textbf{#1}}
\newcommand{\code}[1]{\texttt{\small #1}}
\newcommand{\lecture}[1]{Lecture~#1}

\newtheorem{definition}{Definition}[section]
\newtheorem{proposition}[definition]{Proposition}

% ---- title block -----------------------------------------------------
% \notestitle{19}{Information retrieval: the lexical foundation}{SciFact (BEIR)}
\newcommand{\notestitle}[3]{%
  \thispagestyle{empty}
  \begin{flushleft}
    {\color{aimlmuted}\small\sffamily
     APPLICATIONS OF MACHINE LEARNING \quad$\cdot$\quad
     BSc Mathematics of Artificial Intelligence}\\[2mm]
    {\color{aimlrule}\rule{\textwidth}{0.8pt}}\\[4mm]
    {\color{aimlmuted}\large Lecture #1 \quad$\cdot$\quad Part V \quad$\cdot$\quad
     Extended notes}\\[3mm]
    {\Huge\bfseries\color{aimlprimary}\baselineskip=1.15\baselineskip #2\par}\vspace{4mm}
    {\color{aimlmuted}\large Dataset: #3}\\[3mm]
    {\color{aimlrule}\rule{\textwidth}{0.8pt}}
  \end{flushleft}
  \vspace{2mm}
  \begin{keybox}[title={Why these notes exist}]
    Lectures 19--22 sit outside G\'eron, the course's primary text. For those
    four lectures \emph{these notes are the primary source}, and they are
    examined on exactly the same terms as the textbook parts. Everything in
    the slide deck for this lecture is here, with the arguments written out at
    length rather than compressed onto a slide; the numbers are the ones the
    accompanying notebook prints, and every one of them is a measurement on
    the corpus named above rather than a figure quoted from a paper.
  \end{keybox}
  \vspace{1mm}
  {\small\tableofcontents}
  \vspace{2mm}
  \noindent{\color{aimlrule}\rule{\textwidth}{0.4pt}}
  \vspace{1mm}
}


\begin{document}
\notestitlefor{3}{I}{Classification and its metrics}{MNIST, detecting one digit}{Chapter 3}

\section{The task, and a first detector}

Two lectures on a regression problem ended with a number and an interval around
it. Today the problem is classification, the metric is not RMSE, and the
baseline matters more than it did in either of the previous two lectures.

\subsection{The brief}

A national postal operator sorts handwritten postcodes automatically: a camera
photographs each digit, a classifier reads it, the envelope goes down a chute.
Their audit team has found that one glyph is misread far more often than the
others. They want every scanned digit that is a \textbf{5} pulled off the line
and sent to a human verification desk.

Note what that is not. It is not ``read the digit''. It is \emph{detect one
digit} --- yes or no, one class against the other nine.

\begin{center}
\small
\begin{tabular}{lp{0.55\textwidth}}
\toprule
\textbf{Question} & \textbf{Answer} \\
\midrule
Supervision & supervised --- every image carries a label \\
Task        & binary classification: 5 or not-5 \\
Learning    & batch; the corpus is fixed and fits in memory \\
Assumption  & next year's handwriting looks like this corpus \\
\bottomrule
\end{tabular}
\end{center}

The fourth is the one that will age. Handwriting drifts, cameras are replaced,
and a model trained on one decade's post is not automatically valid on the
next.

\subsection{What makes this a rare-event problem}

The ten digits appear in roughly equal numbers. The \emph{task} does not:
\[ \text{one class against nine} \;\Longrightarrow\;
   \text{about one positive in ten}. \]

\begin{keybox}[title={This is not a fact about MNIST}]
Every binary detector carved out of a multiclass problem inherits this. Detect
one product fault out of forty; one disease out of a hundred; one fraudulent
transaction in ten thousand. The arithmetic is identical and only the ratio
changes.

Nothing about MNIST is imbalanced. Our \emph{question} is.
\end{keybox}

\subsection{The data, and the first default nobody asked for}

MNIST: 70{,}000 images of handwritten digits, each $28\times28$ pixels
flattened to 784 features, each feature one pixel's intensity from 0 (white) to
255 (black). The split is standard --- the first 60{,}000 rows train, the last
10{,}000 test --- and using it is what makes any published MNIST number
comparable to ours.

The labels arrive as one-character \emph{strings}: \code{'5'}, not \code{5}.

\begin{failbox}[title={What forgetting the cast costs}]
\code{y == 5} is then \code{False} for every image, with no error and no
warning, because a string is never equal to an integer. You would build a
detector for a class with no members and score 100\% accuracy on it.

The repair is one line and one assertion:
\code{y = y.astype(np.uint8)}, then
\code{assert y\_train\_5.sum() == 5421}. If the cast was skipped, the assertion
fails at $0 = 5421$ instead of silently training on nothing.
\end{failbox}

\subsection{Metric before model}

Choose what you are optimising before you have anything to optimise ---
otherwise the metric gets chosen after the result, which is how you end up
reporting whichever number happens to look best. And name who chose it: if the
metric arrived by default, from a library or a habit or an assistant, then
nobody chose it.

The obvious metric is \term{accuracy}, the fraction of instances the classifier
gets right. It is between 0 and 1, bigger is better, it needs no threshold and
no scaling, and every stakeholder understands it without being taught.

An \code{SGDClassifier} behind a \code{StandardScaler}, three-fold
cross-validated, scores \textbf{96.90\%} --- folds of 96.83\%, 96.84\% and
97.04\%, agreeing to about two tenths of a point.

\subsection{Before you commit: what does nothing score?}

The cheapest possible detector looks at nothing and answers ``not a 5'' every
time. It is right 54{,}579 times out of 60{,}000: \textbf{90.96\%} accuracy,
with no model, no fit and no features.

\begin{center}
\small
\begin{tabular}{lr}
\toprule
\textbf{Detector} & \textbf{Cross-validated accuracy} \\
\midrule
Never fires --- the anchor      & 90.96\% \\
SGD on raw pixels               & 95.03\% \\
SGD in a pipeline with scaling  & 96.90\% \\
Perfect                         & 100\% \\
\bottomrule
\end{tabular}
\end{center}

Our detector is 5.94 percentage points above doing nothing and 3.10 points
below perfect; of the mistakes the anchor makes we have removed 65.7\%. Those
are the same measurement stated two ways, and you should say which one you
mean.

Read the anchor in two steps. First \emph{what it is}: with no positive
prediction, accuracy is exactly the negative class's base rate, an identity
about the data with no model in it. Only then \emph{what it means}: that
accuracy is available for free, so 90.96\% is zero in the units that matter.
The anchor differs for every task --- 88.76\% for digit 1, because the corpus
holds more 1s than any other digit.

\section{Derivation: imbalance, and the non-monotonicity of precision}

One question, for the whole lecture: doing nothing scores 90.96\% and we score
96.90\%. \emph{What, exactly, did those 5.94 points buy?} The number is
correctly cross-validated, there is no leakage in it, and the three folds agree
to two tenths of a point. Everything about how it was measured is right, which
is what makes the question worth twenty minutes.

\subsection{Notation}

$m$ instances, each with a true label $y^{(i)} \in \{0,1\}$; $P$ positives, $N$
negatives, $P + N = m$; $p = P/m$ the \term{base rate}. A scoring function $s$
and a threshold $t$, predicting positive when $s(x) \ge t$, with counts
$\mathrm{TP}(t)$, $\mathrm{FP}(t)$, $\mathrm{FN}(t)$, $\mathrm{TN}(t)$.

Everything below is arithmetic on those four counts. There is no probability
theory in this derivation and none is needed.

\subsection{Claim 1 --- accuracy is a weighted average of two rates}

\[ \text{accuracy}(t) \;=\;
   \frac{\mathrm{TP}(t) + \mathrm{TN}(t)}{m}. \]

Two of the four counts in the numerator, all four in the denominator --- and
note that the two mistakes enter with \emph{equal weight}. A missed 5 and a
wrongly flagged 3 cost the same, as far as this formula is concerned. Nobody in
the postal operator believes that.

\begin{derivbox}[title={The decomposition}]
Split the four counts by true class, writing
$\text{recall} = \mathrm{TP}/P$ and $\text{specificity} = \mathrm{TN}/N$:
\[
\begin{aligned}
  \text{accuracy}
  &= \frac{\mathrm{TP} + \mathrm{TN}}{m}
   = \frac{P}{m}\cdot\frac{\mathrm{TP}}{P}
   + \frac{N}{m}\cdot\frac{\mathrm{TN}}{N} \\[4pt]
  &= p\,\text{recall} + (1-p)\,\text{specificity}.
\end{aligned}
\]
The classifier that never fires is the special case $\text{recall} = 0$,
$\text{specificity} = 1$, giving
\[ \text{accuracy}(h_0) \;=\; \frac{0 + N}{m} \;=\; 1 - p. \]
\end{derivbox}

An identity, not an estimate: no model, no data, no variance. On our task
$p = 0.09035$, so $\text{accuracy}(h_0) = 0.90965$. That is the 90.96\% we
measured, and it could have been written down without running anything.

\begin{failbox}[title={Corollary --- accuracy is unbounded above by uselessness}]
\[ \lim_{p \to 0^{+}} \text{accuracy}(h_0) \;=\; 1. \]
For any target accuracy $a < 1$ you care to name, there is a detection problem
rare enough that the empty model beats it.

So ``99.2\% accurate'' is not a statement about a classifier. It is a statement
about a classifier \emph{and} a base rate, and quoting it without the second is
quoting half a measurement.
\end{failbox}

\subsubsection{The identity, on our own detector}

\begin{center}
\small
\begin{tabular}{lll}
\toprule
\textbf{Term} & \textbf{Value} & \textbf{Contribution} \\
\midrule
$(1-p)\cdot\text{specificity}$ & $0.90965 \times 0.98860$ & 0.89928 \\
$p\cdot\text{recall}$          & $0.09035 \times 0.77181$ & 0.06973 \\
accuracy                       & ---                      & 0.96902 \\
\bottomrule
\end{tabular}
\end{center}

92.8\% of our headline number is the first row --- being right about the digits
that are \emph{not} 5s. Everything the detector does about the class we were
hired to find is the 0.0697 on the second row.

Stated as a derivative:
\[ \frac{\partial\,\text{accuracy}}{\partial\,\text{recall}} = p,
   \qquad
   \frac{\partial\,\text{accuracy}}{\partial\,\text{specificity}} = 1-p. \]
Accuracy responds to the class you care about in proportion to \emph{how little
of it there is}. You can lose a third of the 5s and pay under three points of
accuracy for it --- and if accuracy is what you optimise, that trade looks
attractive.

\subsection{Claim 2 --- recall is monotone in the threshold and precision is not}

Every scoring classifier is really a family of classifiers, one per threshold:
$\hat y_t(x) = \mathbf{1}[s(x) \ge t]$. Raising $t$ makes it more reluctant to
fire. Nothing about the model changes --- only where we cut.

The flagged set $\{x : s(x) \ge t\}$ shrinks as $t$ rises, and it shrinks by
losing instances, never by gaining any. So for $t' \ge t$, both
$\mathrm{TP}(t') \le \mathrm{TP}(t)$ and $\mathrm{FP}(t') \le \mathrm{FP}(t)$.

\[ \text{recall}(t) = \frac{\mathrm{TP}(t)}{\mathrm{TP}(t) + \mathrm{FN}(t)}
   = \frac{\mathrm{TP}(t)}{P}. \]

The denominator does not depend on $t$ at all: $\mathrm{TP} + \mathrm{FN}$
counts every positive in the data, whatever we predict about it. So recall is a
non-increasing quantity divided by a constant, and recall is non-increasing.
There are no exceptions and no conditions.

Precision has no such guarantee:
\[ \text{precision}(t) =
   \frac{\mathrm{TP}(t)}{\mathrm{TP}(t) + \mathrm{FP}(t)}, \]
where numerator and denominator both move, in the same direction. A ratio of
two decreasing quantities can do anything at all. The denominator is the size
of the set \emph{we chose to flag}: not a property of the data, a property of
our decision.

\begin{derivbox}[title={The one-step lemma}]
Raise $t$ just enough to drop the lowest-scoring flagged instance. Write
$T = \mathrm{TP}$, $F = \mathrm{FP}$, $n = T + F$ before the step, and let
$y \in \{0,1\}$ be the true label of the instance dropped.
\[ \text{precision before} = \frac{T}{n}, \qquad
   \text{precision after}  = \frac{T - y}{n - 1}. \]
Cross-multiplying, the step increases precision exactly when
$n(T-y) > T(n-1)$, that is when
\[ y \;<\; \frac{T}{n} \;=\; \text{precision before}. \]
\end{derivbox}

Since $y$ is 0 or 1, the lemma resolves into two cases:

\begin{center}
\small
\begin{tabular}{llp{0.42\textwidth}}
\toprule
\textbf{Instance dropped} & \textbf{$y$} & \textbf{Effect on precision} \\
\midrule
a false positive & 0 & rises --- always, since precision $> 0$ \\
a true positive  & 1 & falls --- unless precision was already 1 \\
\bottomrule
\end{tabular}
\end{center}

\begin{keybox}[title={The whole result in one sentence}]
Raising the threshold past a positive lowers precision; raising it past a
negative raises precision. The non-monotonicity is not an anomaly --- it is the
generic case. Recall meanwhile falls by $1/P$ in the first case and by nothing
in the second, so both metrics are decided by the same single fact: what the
dropped instance was.
\end{keybox}

\subsubsection{Counted on sixty thousand instances}

Sort every training instance by its out-of-fold score and walk the threshold
down the list one instance at a time --- 59{,}999 steps.

\begin{center}
\small
\begin{tabular}{lr}
\toprule
\textbf{As the threshold rises, one step at a time} & \textbf{Steps} \\
\midrule
the instance dropped is not a 5 --- precision rises & 54{,}579 \\
the instance dropped is a 5 --- precision falls     &  5{,}417 \\
\ldots or stays at 1, because it was already 1      &      3 \\
\bottomrule
\end{tabular}
\end{center}

54{,}579 is the number of non-5s in the training set, and $5{,}417 + 3 =
5{,}420$ is the number of 5s less the one at the very top of the ranking, which
has no step above it. Every step is accounted for by the lemma. This is not a
tendency; it is an exact classification of all 59{,}999 steps.

What is guaranteed, then: recall is non-increasing, always; the flagged count
is non-increasing, always; and precision has no monotone envelope you can rely
on, only a local rule. That is worth remembering because \emph{average
precision} is defined using the maximum precision at or above each recall
level, and that maximum exists for exactly one reason --- to replace a
non-monotone quantity by a monotone one.

\section{What 96.90\% was hiding}

\begin{center}
\small
\begin{tabular}{lrr}
\toprule
\textbf{Detector} & \textbf{Accuracy} & \textbf{Recall} \\
\midrule
Never fires & 90.96\% & 0\% \\
Ours        & 96.90\% & 77.18\% \\
Difference  & 5.94 points & 77.18 points \\
\bottomrule
\end{tabular}
\end{center}

The identity says where the 5.94 points came from. Recall earns $p \times
77.18\% = 6.97$ points; specificity fell from a perfect 100\% to 98.86\%, which
costs $(1-p) \times 1.14\% = 1.04$ points. The three numbers on the page are
each rounded to two decimals, which is why they do not quite close: exactly,
$6.9733 - 1.0370 = 5.9363$ against a measured difference of $5.9370$.

\subsection{The confusion matrix}

\begin{center}
\small
\begin{tabular}{lrr}
\toprule
& \textbf{Predicted not-5} & \textbf{Predicted 5} \\
\midrule
Actually not-5 & 53{,}957 & 622 \\
Actually 5     & 1{,}237  & 4{,}184 \\
\bottomrule
\end{tabular}
\end{center}

Rows are the actual class, columns the predicted class; the diagonal is what
went right. Both mistakes are visible \emph{separately}, and they are not the
same size or the same kind. Accuracy added them together and divided by
60{,}000.

\begin{center}
\small
\begin{tabular}{lrp{0.42\textwidth}}
\toprule
\textbf{Metric} & \textbf{Value} & \textbf{Reads as} \\
\midrule
Precision & 87.06\% & of the digits we flag, 87\% really are 5s \\
Recall    & 77.18\% & of the 5s that exist, we find 77\% \\
Accuracy  & 96.90\% & --- \\
\bottomrule
\end{tabular}
\end{center}

1{,}237 fives went past unflagged: 22.8\% of every 5 in the training set. Both
of those numbers were computable an hour ago. Nothing new was fitted; we asked
the same predictions a different question.

\begin{failbox}[title={Two sentences, both true}]
``The detector is 96.9\% accurate.'' True. Signed off. Deployed.

``The detector misses 22.8\% of the 5s --- on the test shift, about 204
envelopes.'' Also true. Same model, same data, same predictions.

Only one of them is a report the audit team can act on, and it is not the one
the default metric produces.
\end{failbox}

Accuracy was not \emph{wrong}. It correctly answered a question --- what
fraction of all decisions were right? --- that nobody in this brief had asked.
It weights the class we care about by $p = 0.09$, adds two errors the client
prices completely differently, and has a floor of 90.96\% that no model can go
below by being cautious. And we chose it deliberately, before we had anything
to optimise. Choosing early is necessary and not sufficient: only the class
counts would have said it was the wrong metric.

\section{Precision, recall, and the curves}

\subsection{Ask for both numbers, and know what $F_1$ costs}

\[ \text{precision} = \frac{\mathrm{TP}}{\mathrm{TP}+\mathrm{FP}},
   \qquad
   \text{recall} = \frac{\mathrm{TP}}{\mathrm{TP}+\mathrm{FN}}. \]
Recall is also called \term{sensitivity} or the \term{true positive rate};
three names, one quantity, and the third returns on the ROC curve.

When one number is unavoidable, $F_1$ is the harmonic mean:
\[ F_1 = \frac{2}{\frac{1}{\text{precision}} + \frac{1}{\text{recall}}}
       = 2\cdot\frac{\text{precision}\cdot\text{recall}}
                    {\text{precision}+\text{recall}}. \]
Ours is 81.82\%, below both numbers it combines. That is not an accident:

\begin{center}
\small
\begin{tabular}{rrrr}
\toprule
\textbf{Precision} & \textbf{Recall} & \textbf{Arithmetic mean} &
\textbf{$F_1$} \\
\midrule
0.90 & 0.90 & 0.900 & 0.900 \\
0.99 & 0.81 & 0.900 & 0.891 \\
1.00 & 0.02 & 0.510 & 0.039 \\
\bottomrule
\end{tabular}
\end{center}

The harmonic mean is dominated by the smaller of the two. A classifier that
flags three digits a year and is right every time has an arithmetic mean above
one half and an $F_1$ near zero.

\begin{keybox}[title={But that is a choice, not a virtue}]
$F_1$ rewards classifiers with \emph{similar} precision and recall. Our brief
does not ask for similar precision and recall. It asks for the 5s to be found,
within a fixed desk capacity.

A filter for children's videos wants high precision and low recall: rejecting a
good video is annoying, letting a bad one through is not. A shoplifter detector
wants high recall and accepts false alarms: a guard checking one costs a
minute, a missed thief costs the stock. Ours is the second shape, with a cap.
\end{keybox}

\subsection{Where the dial is}

Every scoring classifier computes a number and compares it to a threshold.
\code{predict()} hides the number and hard-codes the threshold at zero. There
is no \code{set\_threshold()} method, and there should not be: you compute the
scores with \code{decision\_function} and compare them yourself.

Choosing 90\% precision on cross-validated training scores gives threshold
20.09, at precision 0.9001 and recall 0.7333 --- ninety per cent precision
costs 3.86 points of recall, from 77.18\% down to 73.33\%.

\begin{failbox}[title={Why the crossing needs a support test}]
We have just spent twenty minutes proving the precision curve is not monotone.
So ``the first index reaching 90\%'' could be a single lucky step, held up by a
handful of flagged digits, that the very next threshold falls back below.

Requiring the crossing to rest on at least 500 flagged digits turns it into a
claim you can defend --- and here it passes unchanged, at threshold 20.09 on
4{,}416 digits across 4{,}411 thresholds.

Aiming higher is nearly free-looking and very expensive. At 99\% precision our
recall falls to 2.12\%: a detector that finds one 5 in fifty. And it fails the
support test --- 99\% precision is reached on only 116 flagged digits, across a
plateau 21 thresholds wide. Which is why 2.12\% is quoted as an illustration of
what high precision costs and never proposed as an operating point. The right
response to a target you cannot meet with support is to \emph{say so}, not to
lower the support test until it passes.
\end{failbox}

A threshold is a hyperparameter, so it is chosen on cross-validated training
scores, never on the test set --- the same error as choosing $\alpha$ that way,
with the same optimistic bias.

\subsection{ROC, and when not to use it}

Same sweep, different axes:
\[ \mathrm{TPR} = \frac{\mathrm{TP}}{P} = \text{recall},
   \qquad
   \mathrm{FPR} = \frac{\mathrm{FP}}{N} = 1 - \text{specificity}. \]
Ours scores ROC AUC $0.9724$, against $0.5$ for a coin and $1.0$ for a perfect
classifier.

Note that \emph{both} denominators are fixed by the data. That is precisely why
the ROC curve behaves so much better than the PR curve --- and why that is a
warning rather than a recommendation.

ROC AUC is the probability that a randomly chosen positive is scored above a
randomly chosen negative. It is a statement about the ranking and nothing else:
it does not depend on the threshold, because there is no threshold in it, and
it does not depend on the base rate, because both axes are normalised within a
class. So it cannot tell you whether the classifier is \emph{usable}. Only how
well it ranks.

\subsubsection{Measure that, rather than taking it on trust}

Take our scores unchanged and thin the positive class. The model, the ranking
and the scores are identical; only the class balance moves.

\begin{center}
\small
\begin{tabular}{lrrr}
\toprule
\textbf{Base rate} & \textbf{Positives} & \textbf{ROC AUC} &
\textbf{Average precision} \\
\midrule
9.04\% --- as measured & 5{,}421 & 0.9724 & 0.8784 \\
2.00\%                 & 1{,}114 & 0.9699 & 0.7203 \\
1.00\%                 &    551 & 0.9709 & 0.5800 \\
\bottomrule
\end{tabular}
\end{center}

ROC AUC moves by 0.0015. Average precision falls by 0.298. Three
indistinguishable ROC curves and three obviously different PR curves, from one
set of scores: \emph{the ROC curve cannot see the problem you have}.

Hence the rule: prefer the precision/recall curve whenever the positive class
is rare, or whenever you care more about false positives than false negatives.
The reason is arithmetic --- FPR divides false positives by $N$, which is huge
when positives are rare, so a large number of false alarms produces a small,
comfortable-looking FPR, while precision divides the same false positives by
the flagged count, which is small.

\subsection{And now a better classifier}

A forest has no \code{decision\_function}. It has \code{predict\_proba}, whose
positive-class column plays exactly the role of the score, and every curve
above works unchanged.

\begin{center}
\small
\begin{tabular}{lrr}
\toprule
\textbf{Metric} & \textbf{SGD} & \textbf{Random forest} \\
\midrule
Accuracy   & 96.90\% & 98.77\% \\
Precision  & 87.06\% & 98.93\% \\
Recall     & 77.18\% & 87.31\% \\
$F_1$      & 81.82\% & 92.76\% \\
ROC AUC    & 0.9724  & 0.9984 \\
Average precision & 0.8784 & 0.9883 \\
Recall available at 90\% precision & 73.33\% & 97.51\% \\
\bottomrule
\end{tabular}
\end{center}

\begin{keybox}[title={Which metric noticed?}]
Accuracy improves by 1.87 points, which reads as polish. Recall improves by
10.13 points --- the thing we care about. Recall at 90\% precision improves by
24.18 points, and those are two completely different systems.

Three descriptions of one comparison. Which one you report decides whether
anyone approves the change. Report the one that matches the brief, and say why
you chose it. \emph{That sentence is the deliverable.}
\end{keybox}

\section{Choosing an operating point}

What the client asked for, in our vocabulary: catch at least 90\% of the 5s on
a shift, and flag no more than 1{,}000 items out of 10{,}000 scanned. Two
constraints, one dial --- and the honest answer may be that no threshold
satisfies both.

The forest at threshold 0.44 was chosen on cross-validated training scores,
where it gives 90.39\% recall at 98.33\% precision. Measured once on the test
shift:

\begin{center}
\small
\begin{tabular}{lrrr}
\toprule
\textbf{Random forest at threshold 0.44} & \textbf{Recall} &
\textbf{Precision} & \textbf{Flagged} \\
\midrule
Cross-validated, on the training set & 90.39\% & 98.33\% & --- \\
Test shift, measured once            & 89.91\% & 98.89\% & 811 \\
\bottomrule
\end{tabular}
\end{center}

Half a point below target. Is that a failure, a bad fold, or noise? It is
noise: the target was hit on a different sample, and a threshold chosen at
exactly 90\% will land either side of it about half the time.

\begin{failbox}[title={Report it, do not tune it away}]
Re-tuning to make the test number reach 90\% is fitting the test set --- the
temptation Lecture 2 told you to refuse, arriving exactly where it was
predicted to.
\end{failbox}

Whichever threshold you propose, you must be able to answer four questions:
which constraint did you satisfy and which did you sacrifice; what does one
extra point of recall cost in items sent to the desk; on what data was the
threshold chosen, and had you seen the test shift when you chose it; and what
happens if the base rate changes next quarter.

\begin{keybox}[title={The general rule}]
Never optimise one metric of a pair --- precision and recall, bias and
variance, latency and accuracy, memory and throughput. Ask for one and you will
get it, at the price of the other, and the price will not be in the reply. The
counter-measure is not cleverness; it is stating the second quantity as a
constraint in the same sentence as the first.
\end{keybox}

\section{Beyond binary}

Everything above was one class against the rest. Three more shapes reuse the
same confusion matrix.

\begin{center}
\small
\begin{tabular}{llp{0.34\textwidth}}
\toprule
\textbf{Shape} & \textbf{Output per instance} & \textbf{Example} \\
\midrule
Binary       & one label from two            & is it a 5? \\
Multiclass   & one label from many           & which digit is it? \\
Multilabel   & several labels at once        & is it large? is it odd? \\
Multioutput  & several labels, each multiclass & the denoised image, pixel by
                                                pixel \\
\bottomrule
\end{tabular}
\end{center}

Multiclass is built from binary parts. \term{One-versus-rest} trains ten
detectors like ours and takes the highest score. \term{One-versus-one} trains
one detector per pair --- $10\times9/2 = 45$ fits, each on a much smaller
training set --- and takes the most-won duel. Scikit-learn picks one for you
automatically and does not mention which: another default nobody asked for, and
it matters, because for algorithms that scale badly with training-set size
forty-five small fits beat ten large ones.

For multilabel targets, metrics are computed per label and averaged, and the
averaging argument is a third default that changes the number without changing
the model: \code{average="macro"} weights every label equally,
\code{average="weighted"} weights by support, and on imbalanced labels the two
can disagree substantially. Name the metric, and name the average.

\section{The notebook}

\begin{trybox}[title={What to change}]
Run it once, then re-run the base-rate experiment with the positive class
thinned to 0.5\% instead of 1\%. ROC AUC will barely move again; watch what
average precision does. Two lines, and it is the whole argument for preferring
the PR curve on rare events --- measured by you rather than asserted here.
\end{trybox}

\section{Exercises}

\begin{enumerate}[leftmargin=1.6em,itemsep=4pt]
  \item A colleague reports 99.2\% accuracy on a rare defect, correctly
        cross-validated. State one further quantity you would need, and why.
        \hfill \scope{[4]}
  \item A ranking classifier is evaluated on 10 instances, of which 4 are
        positive. Sorted by score, highest first, the true labels are
        $1,1,0,1,0,0,1,0,0,0$. Compute precision and recall when the top 4 are
        flagged and when the top 5 are flagged; then state, using the one-step
        lemma, whether precision rises or falls when the threshold is raised
        from ``top 5'' to ``top 4'', and why. \hfill \scope{[3]}
  \item Explain why $F_1$ would be the wrong summary for the postal operator's
        brief, and name the pair of numbers you would report instead.
        \hfill \scope{[4]}
  \item A model's ROC AUC is unchanged when the positive class is thinned from
        9\% to 1\%, while its average precision falls by 0.30. Explain both
        observations in terms of the denominators of the two curves' axes.
        \hfill \scope{[5]}
  \item You are asked for a threshold achieving 99\% precision. The first
        crossing rests on 116 flagged instances across a plateau 21 thresholds
        wide. State what you would report, and why lowering the support test
        until the target passes is not an option. \hfill \scope{[4]}
\end{enumerate}

\section*{Summary}
\addcontentsline{toc}{section}{Summary}

Accuracy decomposes as $p\,\text{recall} + (1-p)\,\text{specificity}$, so a
classifier that never fires scores exactly $1-p$ and
$\partial\,\text{accuracy}/\partial\,\text{recall} = p$. On this task that
anchor is 90.96\%, our detector scored 96.90\%, and 92.8\% of the headline was
being right about digits that are not 5s.

Recall is non-increasing in the threshold because its denominator is $P$, which
the threshold cannot touch. Precision's denominator is the flagged count, which
the threshold chooses, so one step raises precision exactly when the dropped
instance is a negative --- counted on our own ranking as 54{,}579 rises,
5{,}417 falls and 3 flat, every one of 59{,}999 steps accounted for.

Asking the same predictions a different question gave precision 87.06\% and
recall 77.18\%: 1{,}237 fives past the desk, 22.8\% of them. ROC AUC moved by
0.0015 as the base rate fell from 9\% to 1\% while average precision fell by
0.298, which is why the PR curve is the one to read on a rare event. And the
forest improved accuracy by 1.87 points, recall by 10.13, and recall at 90\%
precision by 24.18 --- three descriptions of one comparison, of which only the
last matches the brief.

\end{document}
